\chapter{Bringup 与配置系统}
\label{ch:bringup}

\section{启动组织}
启动文件只负责编排节点、声明参数、选择 profile、建立 remap 和处理关键进程退出；算法状态和业务逻辑留在节点实现中。真机启动由 \pkg{uv_bringup} 管理，仿真启动由 \pkg{uv_sim_bringup} 管理。

\begin{figure}[H]
  \centering
  \begin{tikzpicture}[
    box/.style={draw=YoulongBlue, rounded corners, fill=YoulongLight,
      minimum width=34mm, minimum height=8mm, align=center},
    arr/.style={-{Latex[length=2mm]}, thick, color=YoulongBlue},
    node distance=7mm
  ]
    \node[box] (real) {\code{real.launch.py}};
    \node[box, below left=of real] (desc) {真实 URDF};
    \node[box, below=of real] (hw) {\pkg{uv_hm}};
    \node[box, below right=of real] (loc) {\pkg{uv_localization}};
    \node[box, below=of hw] (ctrl) {控制 / 感知 / 规划 / 任务};
    \draw[arr] (real) -- (desc);
    \draw[arr] (real) -- (hw);
    \draw[arr] (real) -- (loc);
    \draw[arr] (desc) |- (ctrl);
    \draw[arr] (hw) -- (ctrl);
    \draw[arr] (loc) -- (ctrl);
  \end{tikzpicture}
  \caption{真机 bringup 组成}
  \label{fig:real-bringup}
\end{figure}

仿真入口 \code{sim.launch.py} 额外启动仿真 URDF、Stonefish、\pkg{uv_sim_bridge}、可选退化节点、评测和观测工具；HIL 入口再启动 micro-ROS Agent 并将 \code{hil_mode} 传给 bridge。

\section{Profile 机制}
profile 是标准 ROS 2 YAML 参数文件，用于按模式选择保守或开发参数。当前可见 profile 包括：
\begin{itemize}
  \item real：\code{real_default}、\code{real_safe}；
  \item sim：\code{sim_dev}、\code{sim_ci}；公开 \pkg{uv_sim} 入口另提供
        \code{guoshui}、\path{guoshui_cruise}、\path{guoshui_cruise_seeded}、
        \code{sauvc_finals} 和 \path{sauvc_qualification} world 预设；
  \item hil：\code{hil_lab}。
\end{itemize}
启动文件在进入子 launch 前解析 profile 路径并校验 profile 与运行模式匹配，避免将仿真参数误用于真机。

\section{运行预设}
\code{profile} 与 \code{preset} 是两层不同的配置。\code{profile} 选择传给节点的 ROS 2 参数文件；\code{preset} 则从独立的 launch YAML 读取一组启动参数默认值。显式写在命令行上的 launch 参数优先于预设值；不指定 \code{preset} 或选择 \code{default} 时保留入口原有默认行为。

预设按运行时和所属包分别维护：真机位于 \pathref{workspace_auv/src/uv_bringup/config/launch_profiles/real/}；SIL 与 HIL 分别位于 \pathref{workspace_sim/src/uv_sim_bringup/config/launch_profiles/sim/} 和 \pathref{workspace_sim/src/uv_sim_bringup/config/launch_profiles/hil/}。每个目录提供 \code{default.yaml}、\code{record.yaml}、\code{debug.yaml} 和 \code{task.yaml}。

\begin{itemize}
  \item \code{record} 启动 \pkg{uv_record} 的 \code{raw} 图像模式，从 Iceoryx2 相机源帧保存无损 PNG，适合后续整理数据集；同时保存 rosbag 和 session 日志，不要求启动 go2rtc。
  \item \code{debug} 默认启用 AI，并录制 go2rtc 未标注相机流（JPEG）及 rosbag、session 日志，便于复现运行问题。
  \item \code{task} 启用 AI、运动、导航和任务执行器，并组合 debug 录制设置。任务执行器会自动加载并开始执行默认的 \code{robocup_26.yaml}；需要其他任务时可显式传入 \code{mission_file}。
\end{itemize}

\Needspace{12\baselineskip}
\section{常用启动命令}
\begin{lstlisting}[language=bash,caption={标准运行入口}]
# Real
cd workspace_auv
colcon build --symlink-install
source install/setup.bash
ros2 launch uv_bringup real.launch.py profile:=real_default

# SIL
cd workspace_sim
colcon build
source install/setup.bash
ros2 launch uv_sim_bringup sim.launch.py profile:=sim_dev

# 公开场景预设；显式 world:= 会覆盖预设 world
ros2 launch uv_sim sim.launch.py profile:=sauvc_finals

# HIL
ros2 launch uv_sim_bringup hil.launch.py profile:=hil_lab \
  serial_dev:=/dev/ttyUSB0 serial_baud:=921600
\end{lstlisting}

\begin{lstlisting}[language=bash,caption={使用运行预设的一键启动入口}]
# 数据集采集：uv_record raw 源帧
ros2 launch uv_bringup real.launch.py preset:=record
ros2 launch uv_sim_bringup sim.launch.py preset:=record
ros2 launch uv_sim_bringup hil.launch.py preset:=record

# 调试：AI + 未标注 go2rtc 视频 + rosbag/session 日志
ros2 launch uv_bringup real.launch.py preset:=debug
ros2 launch uv_sim_bringup sim.launch.py preset:=debug
ros2 launch uv_sim_bringup hil.launch.py preset:=debug

# 一键运行默认 Mission；也可附加 mission_file:=/path/to/mission.yaml
ros2 launch uv_bringup real.launch.py preset:=task
ros2 launch uv_sim_bringup sim.launch.py preset:=task
ros2 launch uv_sim_bringup hil.launch.py preset:=task
\end{lstlisting}

\begin{warningbox}
\code{task} 是会自动执行 Mission 的运行预设，不只是启动一个空闲的任务节点。真机入口同时启用运动控制；使用前应确认车辆处于允许运动的安全环境，并检查实际传入的 \code{mission_file}。
\end{warningbox}

\section{启动顺序与就绪}
典型依赖顺序是：描述/固定 TF、硬件或仿真输入、定位、控制、相机/感知、规划、任务。常规启动默认不执行任务；选择 \code{preset:=task} 时会显式启用任务执行器并自动开始默认 Mission。\pkg{uv_bringup} 提供 readiness 检查，可按 backend、control、sensors、perception 阶段确认 Topic 和节点已达到运行条件。

\section{参数事实来源}
启动参数的默认值在 launch 中定义，模式参数在 profile 中覆盖，任务参数在 Mission/Task YAML 中维护，相机参数在相机 profile 和标定文件中维护。参数优先级和最终生效值应通过启动日志或 ROS 参数查询确认；LaTeX 不复制完整参数树，附录只列出入口和语义。
